%% titre: Introduction à la programmation linéaire
%% niveau: Licence 3
%% matiere: Recherche opérationnelle
%% annee: 2025-2026
%% ordre: 1
%% description: modélisation, forme standard et méthode graphique.

\section{Modélisation}

Un \emph{programme linéaire} consiste à optimiser une fonction linéaire sous des contraintes linéaires. Soit $c\in\R^n$, $A\in\mathcal{M}_{m,n}(\R)$ et $b\in\R^m$ ; on considère
\[
  (P)\quad \min_{x\in\R^n}\ c^\top x \quad\text{sous}\quad Ax\le b,\ x\ge 0 .
\]

\begin{definition}[Domaine admissible]
On appelle \textbf{domaine admissible} l'ensemble
$$\mathcal{D}=\{x\in\R^n \mid Ax\le b,\ x\ge 0\}.$$
\end{definition}

\begin{theoreme}[Théorème fondamental]
Si $(P)$ admet une solution optimale, alors il en existe une qui est un sommet de $\mathcal{D}$.
\end{theoreme}

\begin{proof}
Le domaine $\mathcal{D}$ est un polyèdre convexe. Une fonction linéaire atteint son minimum sur un polyèdre en au moins un point extrémal, c'est-à-dire en un sommet.
\end{proof}

\begin{remarque}
Les contraintes d'égalité $Ax=b$ se ramènent à deux inégalités $Ax\le b$ et $-Ax\le -b$.
\end{remarque}

\section{Forme standard}

\subsection{Variables d'écart}

Pour passer à la forme standard, on introduit des variables d'écart $s_i\ge 0$ :
\begin{align*}
  a_{i1}x_1+\dots+a_{in}x_n \le b_i
  &\iff a_{i1}x_1+\dots+a_{in}x_n + s_i = b_i,
\end{align*}
et l'on obtient $\min\ c^\top x$ sous $Ax+s=b$, $x,s\ge 0$.

\subsection{Un exemple}

\begin{exemple}
Un atelier fabrique deux produits $P_1$ et $P_2$ :
\begin{center}
\begin{tabular}{|l|c|c|c|}
\hline
 & $P_1$ & $P_2$ & Disponible \\
\hline
Heures de machine & 2 & 1 & 100 \\
Heures de main-d'\oe uvre & 1 & 1 & 80 \\
\hline
Profit unitaire (\texteuro) & 30 & 20 & \\
\hline
\end{tabular}
\end{center}
Le programme s'écrit
\[
  \max\ 30x_1+20x_2 \quad\text{sous}\quad
  \begin{cases}
    2x_1+x_2\le 100\\
    x_1+x_2\le 80\\
    x_1,x_2\ge 0.
  \end{cases}
\]
\end{exemple}

\section{Méthode graphique}

Les étapes sont les suivantes :
\begin{enumerate}[label=\arabic*)]
  \item tracer le domaine admissible ;
  \item tracer une droite de niveau $c^\top x = \lambda$ ;
  \item la déplacer dans la direction d'amélioration jusqu'à la dernière intersection avec $\mathcal{D}$.
\end{enumerate}

\begin{exercice}[Application]
Résoudre graphiquement l'exemple ci-dessus.
\begin{correction}
Les sommets du domaine sont $(0,0)$, $(50,0)$, $(20,60)$ et $(0,80)$. La fonction objectif vaut respectivement $0$, $1500$, $1800$ et $1600$ : l'optimum est atteint en $(20,60)$, pour un profit de $1800$.
\end{correction}
\end{exercice}
